\documentclass[12pt]{article}
\usepackage{amsmath,amsfonts,amsthm}
\newtheorem{theorem}{Theorem}
\usepackage{fullpage}
\usepackage{natbib}
\usepackage{algpseudocode}
\usepackage{algorithm}

\begin{document}

\author{Toby Hocking}

\title{Linear time functional dynamic programming for un-regularized isotonic regression}

\maketitle

\section{Introduction}

Isotonic regression is a statistical model which fits a non-decreasing function to a sequence of $n$ data.
Previous algorithms for computing this model include the Pool Adjacent Violators Algorithm (PAVA), which can be implemented in linear time, $O(n)$ \citep{Busing2022}.
We propose a new dynamic programming algorithm, and demonstrate the similarities and differences with PAVA.

\section{Problem}

\subsection{Data}

We assume there are $n$ data, $y_1,\dots,y_n\in \mathbb R$.
In some contexts, there are also positive weights $w_1,\dots,w_n\in \mathbb R_+$.
When no weights are given, we take $w_i=1$ for all $i$.

\subsection{Isotonic regression problem}

The weighted isotonic regression problem is 
\begin{eqnarray}
    \label{eq:isotonic-problem}
\hat C_n &= & \min_{x_1,\, \dots,\, x_{n}}  \sum_{i=1}^n w_i(x_i-y_i)^2 \text{ subject to } \forall i\in\{2, \dots, n\},\, x_{i-1} \leq x_{i}.
\end{eqnarray}
The $\hat C_n\in R$ is the optimal cost, including all $n$ data points.

\section{Algorithms}

\subsection{Proposed dynamic programming algorithm}
We propose a new algorithm based on previous work from the literature on optimal change-point detection \citep{Hocking2020jmlr}.
The idea is to use dynamic programming to recursively compute cost functions,
\begin{eqnarray}
    \label{eq:Cn_mu}
  C_n(x_n) &= & \min_{x_1,\, \dots,\, x_{n-1}}  \sum_{i=1}^n w_i(x_i-y_i)^2 \text{ subject to } \forall i\in\{2, \dots, n\},\, x_{i-1} \leq x_{i}.\\
           &=& w_n(x_n-y_n)^2 + \min_{x_{n-1}} C_{n-1}(x_{n-1}) \text{ subject to } x_{n-1}\leq x_n.\\
  &=& w_n(x_n-y_n)^2 + C_{n-1}^\leq(x_{n}),
\end{eqnarray}
where $C_{n-1}^\leq : \mathbb R \rightarrow \mathbb R$ is a function that we call the min-less operator of the optimal cost function up to $n-1$.
Dynamic programming means that we have a simple update rule for computing the optimal cost up to $n$ data points, using the optimal cost up to $n-1$ data points.
Functional representation means that the equations above use optimal cost functions rather than values.

\subsection{Functional representation}

Whereas previous functional dynamic programming algorithms have used either a linked list or binary tree data structure, the proposed algorithm uses arrays, which is a much simpler design that is only possible because of the problem structure.
We will show that $C_i(x)$ is a convex function, defined in terms of pieces, each piece corresponding to a cluster of data with the same mean, and the min is always on the last piece (TODO examples and proof).

Each iteration $i\in\{1,\dots,n-1\}$ of dynamic programming computes $F_{i+1}$, a cost function with $K_i+1\in\{1,\dots,i+1\}$ pieces.
Let $M_{i,0}=-\infty < M_{i,1} < \cdots < M_{i, K_i} < M_{i, K_i+1}=\infty$ be bounds on the function pieces, at iteration $i$ of the dynamic programming algorithm.
Let $f_{i,k}$ be the quadratic function piece $k$ at iteration $i$.
Then we have
\begin{eqnarray}
  \label{eq:update}
  F_{i+1}(x) &=& (x-y_{i+1})^2 + F_i^\leq(x) \\
  &=&
      \begin{cases}
        f_{i,1}(x) &\text{ if } x\in(-\infty=M_{i,0},\, M_{i,1}],\\
         &\vdots\\
        f_{i,K_i+1}(x) &\text{ if } x\in[M_{i,K_i},\, M_{i,K_i+1}=\infty).
      \end{cases}
\end{eqnarray}
For each piece $k\in\{1,\dots,K_i+1\}$, there is a corresponding convex quadratic function.
Let $\hat N_{i,k}$ be the number of data at iteration $i$ in piece $k$, and let $\hat T_{i,k}$ be the total data values at iteration $i$ in piece $k$.
Then we have
\begin{equation}
  \label{eq:fik}
  f_{i,k}(x) = x^2\hat N_{i,k}  -2 x\hat T_{i,k}  + \hat c_{i,k},
\end{equation}
where $\hat N_{i,k}, \hat T_{i,k}$ are coefficients that can be efficiently computed during the algorithm, if necessary.
It is not necessary to store/compute $\hat c_{i,k}$, because it does not appear in the derivative,
\begin{equation}
  \label{eq:fik_deriv}
  f'_{i,k}(x) = 2x\hat N_{i,k}  -2\hat T_{i,k} = 0 \Rightarrow x=\hat T_{i,k}/\hat N_{i,k}.
\end{equation}

\subsection{Details of dynamic programming update rules}

For each iteration $i\in\{1, \dots, n\}$, the dynamic programming stores the current number of clusters $K_i\in\{1,\dots,i\}$, and arrays for cluster size and mean $N_{i,k},M_{i,k}$, for $k\in\{1,\dots,K_i\}$.

The initialization is to set the number of clusters to one $K_1=1$, set the cluster size to one, $N_{1,1}=1$, and set the cluster mean to the first data point, $M_{1,1}=y_1$.
This is a representation of the first cost function, $F_1^\leq$.

For each iteration $i\in\{1, \dots, n-1\}$, the dynamic programming recursively computes $F_{i+1}^\leq$, via
\begin{equation}
  \label{eq:Kip1}
  K_{i+1} = \max\{
  k\in\{K_i+1,\dots,1\}
  \mid
  f'_{i,k}( M_{i,k-1} ) < 0
  \Rightarrow
  \hat T_{i,k}/\hat N_{i,k} > M_{i,k-1}
  \}.
\end{equation}
The condition $  f'_{i,k}( M_{i,k-1} ) < 0$ checks to see if the minimum occurs in the interval $(M_{i,k-1},M_{i,k})$.
The max starts by checking $k=K_i+1$, then decrements $k$ if the condition is not satisfied. 

The coefficients of $f_{i,k}$ are computed via a cumulative sum (starting at $k=K_i+1$ and going down),
\begin{eqnarray}
  \label{eq:cum}
  \hat N_{i, k} &=& 
  \begin{cases}
    1 & \text{ if } k=K_i+1.\\
    \hat N_{i, k+1} + N_{i,k} = 1+\sum_{j=k}^{K_i} N_{i,j} & \text{ if } k\in\{K_i,\dots,1\}.
  \end{cases}\\
  \hat T_{i, k} &=& 
  \begin{cases}
    y_{i+1} & \text{ if } k=K_i+1.\\
    \hat T_{i, k+1} + N_{i,k} M_{i,k} = y_{i+1}+\sum_{j=k}^{K_i} N_{i,j} M_{i,j} & \text{ if } k\in\{K_i,\dots,1\}.
  \end{cases}
\end{eqnarray}
At step $i+1$, the new array values are
\begin{eqnarray}
  \label{eq:Nip1}
  N_{i+1, k} &=& 
  \begin{cases}
    N_{i,k} & \text{ for } k\in\{1,\dots,K_{i+1}-1\},\\
    \hat N_{i,K_{i+1}} & \text{ for } k=K_{i+1}.
  \end{cases}  \\
  \label{eq:Mip1}
  M_{i+1, k} &=& 
  \begin{cases}
    M_{i,k} & \text{ for } k\in\{1,\dots,K_{i+1}-1\},\\
    \hat T_{i,K_{i+1}}/\hat N_{i,K_{i+1}} & \text{ for } k=K_{i+1}.
  \end{cases}  
\end{eqnarray}
Note that values for all $k$ except the largest ($k=K_{i+1}$) are the same between iterations $i$ and $i+1$.
We therefore can avoid allocating new arrays, and simply write the new values in position $K_{i+1}$ of the existing arrays.

\begin{theorem}
  The recursively computed values $K_n,N_{n,k},M_{n,k}$ define a function $F_n^\leq$ which is equal to the optimal cost, $C_{n}^\leq$.
\end{theorem}
\begin{proof}
TODO
\end{proof}

\subsection{Pseudocode}
\begin{algorithm}
\begin{algorithmic}
  \Require data $y_1,\dots, y_n\in\mathbb R$.
  \State Allocate number of clusters $K\gets 0$, cluster sizes $N\in\mathbb Z^n$ and means $M\in\mathbb R^n$.
  \For{$i=1$ to $n$}
    \State $\texttt{samples} \gets 1$, $\texttt{total} \gets y_i$.
    \While{$K>0$ and $\texttt{total}/\texttt{samples}\leq M_K$}
      \State $\texttt{samples } += N_K$, $\texttt{total } += N_K M_K$, $K--$.
    \EndWhile
    \State $K++$, $N_K\gets\texttt{samples}$, $M_K\gets\texttt{total}/\texttt{samples}$.
  \EndFor
  \State \Return $K,N,M$
\end{algorithmic}
  \caption{\label{alg:iso-dp}Functional dynamic programming for isotonic regression}
\end{algorithm}

Algorithm~\ref{alg:iso-dp} summarizes the solution method for the simple case without weights ($w_i=1$ for all $i$).
It outputs the optimal means $M$ along with cluster sizes $N$, which is a run-length encoding of the solution, for a number of clusters $K\in\{1,\dots,n\}$.
To convert into the $n$-dimensional real-valued solution (an optimal mean value for each input data point), first set all $N$ values after index $K$ to zero, then $x\gets \texttt{rep}(M,N)$ (R code for repeating each $M_k$ value $N_k$ times).

\subsection{Comparison with PAVA}
TODO

\subsection{Extensions for bounds and weights}
TODO

\bibliographystyle{abbrvnat}
\bibliography{refs}

\end{document}
